% reported checkpoint appendix draft generated from the frozen run manifest and implementation.
% Replace model-version wording only if the camera-ready paper anonymizes vendors.

\section{Framework Details}
\label{app:framework_detail}

\input{tex_appendix/algorithm}

\subsection{Rubric Construction and Validation}

For each completed Doctor--Agent trajectory, we first generated an
answer-neutral, post-hoc behavioral reflection.  The reflection linked three
observable elements when available: (i) the physician's explicit request,
objection, uncertainty, or concern; (ii) the handling of that issue in one
specific Agent reply; and (iii) the physician's subsequent response indicating
that the issue was addressed, remained unresolved, or that the physician
maintained their position.  The reflection ended by selecting exactly one
behavioral locus, denoted by an Agent reply identifier.  The reflection model
did not receive the question, answer options, gold answer, final decision, or
trajectory outcome.  In the frozen reported run, this role used Claude Haiku 4.5
with the \texttt{behavior\_observer\_v5} protocol.  We refer to this signal as
partner-side reflection: it describes the physician's observed interaction
with the Agent, but it is generated by a separate observer rather than by the
local Doctor model itself.

The rubric generator received the completed trajectory, its success/failure
stratum, the answer-neutral reflection, and the selected behavioral locus.  It
produced one behavioral construct with (i) a physician-side applicability
condition, (ii) a condition on the Agent's visible state before the target
reply, (iii) a required response behavior, and (iv) ordinal descriptions for
scores 0, 1, and 2.  Applicability was restricted to information visible
strictly before the scored reply.  Later turns, the physician's final decision,
the outcome label, and hidden state could not establish applicability.  The
same applicability statements were repeated across all three score levels so
that the levels differed only in the observable response behavior.

For a successful source trajectory, score 2 represented the helpful behavior
observed at the selected locus, score 1 represented partial performance of that
same behavior, and score 0 removed only that behavior.  For a failed source
trajectory, score 0 represented the observed handling problem or omission,
score 1 represented a partial correction, and score 2 corrected that same
problem.  Scores were not definitions of medical correctness: a correct final
answer did not imply score 2, and an incorrect final answer did not imply score
0.  Public criteria were answer-neutral and excluded option labels, answer
text, and case-specific medical entities.

To reduce duplicate criteria, the generator could assign a later trajectory to
an existing rubric only when the interaction construct, visible applicability,
required behavior, success/failure direction, and all three score semantics
transferred without reinterpretation or a new premise.  Matching was performed
within the same case and source-outcome stratum in the frozen implementation;
topic similarity alone was insufficient.

Each newly created rubric underwent deterministic schema and leakage checks,
LLM-based checks for answer presupposition and unsupported domain premises, and
a source-locus consistency check.  A surviving rubric was then subjected to an
intervention test.  For a failed source, the intervention requested score-2
handling; for a successful source, it selectively removed the score-2 behavior
to request score-0 handling.  The policy generated one intervened Agent reply,
after which the Doctor produced one fresh final decision.  Up to three reruns
were used, with early resolution after two agreeing outcomes.  A candidate was
retained when the desired outcome change occurred in at least two of three
reruns on a valid sampled witness.  The bounded-witness implementation sampled
at most ten same-rubric trajectories in deterministic random order and retained
the rubric after the first passing witness.  Applicability failures consumed a
sampled slot but did not receive a behavioral score.  Intervention adherence
was not an additional reward or retention criterion.

\subsection{Reward Assignment and Policy Optimization}

Only retained rubrics contributed to response scoring.  Applicability was
evaluated from the strict pre-response state and was shared by sibling replies
from the same state.  Conditional on applicability, each reply was independently
assigned a level $\ell_{i,r}\in\{0,1,2\}$ for rubric $r$.  We normalized this
level as $s_{i,r}=\ell_{i,r}/2$ and averaged equally over the applicable retained
rubrics,
\begin{equation}
R_A(i)=\frac{1}{|\mathcal{A}_i|}\sum_{r\in\mathcal{A}_i}s_{i,r},
\end{equation}
where $\mathcal{A}_i$ is the set of applicable retained rubrics for response
$i$.  The reported checkpoint condition was rubric-only, so the training reward was
$R(i)=R_A(i)$ and the final task outcome had weight zero.  When
$\mathcal{A}_i=\varnothing$, the response remained unassessed and was excluded
from the rubric-only update rather than being converted to reward zero.

The optimization unit was a same-state sibling group.  With the mean-centering
mode used in reported checkpoint, the advantage was
\begin{equation}
\hat A_i = R(i)-\frac{1}{|G|}\sum_{j\in G}R(j),
\end{equation}
for assessed responses in sibling group $G$.  Groups with fewer than two
assessed responses or with all-zero advantages were skipped.  Thus, reward
comparisons were made only among responses sampled from the same dialogue state.

Let $\pi_{\mathrm{old}}$ be the behavior policy that sampled a response,
$\pi_\theta$ the trainable policy, and $\pi_{\mathrm{ref}}$ the frozen base
model with the LoRA adapter disabled.  For completion token $t$, we used
$r_{i,t}=\pi_\theta(y_{i,t})/\pi_{\mathrm{old}}(y_{i,t})$ and optimized the
clipped GRPO objective
\begin{align}
\mathcal{L}_{i} ={}& -\frac{1}{T_i}\sum_t
\min\!\left(r_{i,t}\hat A_i,
\operatorname{clip}(r_{i,t},1-\epsilon,1+\epsilon)\hat A_i\right) \\
&+\frac{\beta}{T_i}\sum_t
\left[\exp(\Delta_{i,t})-\Delta_{i,t}-1\right],
\end{align}
where $\Delta_{i,t}=\log\pi_{\mathrm{ref}}(y_{i,t})-
\log\pi_\theta(y_{i,t})$, $\epsilon=0.2$, and $\beta=0.04$.
Teacher-forced log probabilities were computed over every sampled Agent token.
Under the precommit protocol, this included both the private pre-assessment and
the visible reply, while Doctor tokens were never optimized.  The implementation
also required the pre-update policy ratio to agree with the recorded behavior
policy within $10^{-3}$.

\section{Implementation Details}
\label{app:implementation_details}

\subsection{Training Configuration}

The Agent was Qwen3-8B loaded with Hugging Face Transformers at revision
\texttt{b968826d9c46dd6066d109eabc6255188de91218}.  We trained a rank-16
RS-LoRA adapter with $\alpha=32$, zero LoRA dropout, and target modules
\texttt{q\_proj}, \texttt{k\_proj}, \texttt{v\_proj}, and
\texttt{o\_proj}.  The learning rate was $10^{-5}$, gradient norm was clipped
to 1.0, and token log probabilities were projected in chunks of 16 tokens.

Training used five NVIDIA A100 80GB GPUs. Five cases formed one synchronization
batch, with one case initially assigned per GPU and the optimizer update
performed on GPU 4 after the batch barrier. The frozen training order and its
SHA-256 digest are included in the artifact.

\subsection{Clinical Collaboration Environment}

The Doctor was Gemma 3 12B Instruction Tuned served locally by Ollama 0.12.5 as
\texttt{gemma3:12b-it-q8\_0} (Q8\_0 quantization; recorded model digest
\texttt{997a7c2c09753e9011ec18a3fcb498e27c0ba330599dfb26a2e05b58cfbdc1bc}).
One Doctor server was pinned to each physical GPU.  Doctor discussion responses
used temperature 0.4, top-$p$ 1.0, and a 256-token limit; final decisions used a
label-only protocol with a 32-token limit.

The Agent used temperature 0.7, top-$p$ 0.9, and a maximum completion length of
400 tokens.  At every decision point it first produced a private pre-assessment
without access to the preceding discussion and then generated its visible reply
conditioned on both that pre-assessment and the collaboration history.  The
rollout was a full response-level tree with branching factor four and at most
four Agent turns per branch.  Adaptive stopping was permitted only after at
least two Agent turns.  A single immutable Doctor opening was shared across
conditions for a case, while subsequent Doctor turns, Agent actions, final
decisions, and rewards remained condition-specific.

\subsection{Rubric Generation and Scoring}

The partner-side reflection used Claude Haiku 4.5 with a 512-token output
limit.  Gemini 3.8 Flash generated rubrics with temperature 0.7, a 4096-token
limit, low reasoning effort, and strict JSON-schema output.  GPT-5.6 Luna
performed the rubric Gate and two-stage scoring.  The first scoring stage
determined applicability without seeing the target response; the second stage
scored each applicable response.  No-match results remained unscored.  The
prompt-search reflection and refinement roles used Claude Sonnet 5.5.  All
prompt files and model-role identifiers are preserved in the run manifest with
their SHA-256 hashes.

Prompt optimization was evaluated once per five-case batch.  One candidate was
generated at refinement depth one and compared with its parent on identical
frozen trees.  Prompt fitness used the product of the mean sampled witness
effects for success-source and failure-source rubrics.  Ties or a missing
direction retained the parent.  Training used the deduplicated union of retained
rubrics from both compared prompts.  These prompt-search reruns and the
rubric-retention decision reused recorded witness results but remained distinct
from response reward assignment.

\subsection{Evaluation and Baseline Details}

At reported checkpoint, rollout and optimization were paused and the frozen LoRA checkpoint was
evaluated on all 622 cases in the fixed CasiMedicos English evaluation set
(SHA-256
\texttt{e42077d006c4e8227564f99b06050b66e7df9989f9b2b17f5fc5d9dd083feb56}).
The collaboration protocol and Doctor configuration matched training.  The
physician's label-only final decision was parsed as the prediction.  All 622
outputs produced one valid option; 442 were correct and 180 were incorrect,
giving an accuracy of $71.0609\%$.

The released reported checkpoint artifact contains the complete per-case predictions and
correctness values.  A paired Fix/Break comparison requires per-case predictions
from the corresponding vanilla baseline on the identical 622-case ordering.
Those baseline predictions are not part of the frozen reported checkpoint artifact described
here; therefore baseline accuracy and Fix/Break counts should be inserted only
after linking the separately preserved baseline result by case ID.  Fix denotes
baseline-incorrect to collaborative-correct, and Break denotes
baseline-correct to collaborative-incorrect.

